% =====================================================================
\section[Preliminaries and the noncommutative matrix-ring framework]%
  {Preliminaries and the noncommutative\\ matrix-ring framework}
\label{sec:prelim}
% =====================================================================

\subsection{Standing hypotheses and notation}

Throughout, $p$ is a prime, $q=p^r$, and $t,n\ge1$ are integers. We set
\begin{equation}\label{eq:rings}
  A_t=\F_q[u]/\langle u^t\rangle,\qquad
  R=M_2(A_t),\qquad
  \Bq=A_t[X]/\langle X^n-1\rangle,\qquad
  \Rn=R[X]/\langle X^n-1\rangle ,
\end{equation}
and write $\Nlen=4tn$. Elements of $A_t$ are written $\sum_{a=0}^{t-1}c_au^a$.
All dimensions and all cardinalities of codes are over $\F_q$ unless a different
field is named explicitly; this is maintained without exception, and the only
place where $\F_p$ enters is Theorem~\ref{thm:fpbasis}.

\begin{definition}[cyclic code]\label{def:code}
A \emph{cyclic code of length $n$ over $R$} is a right ideal of $\Rn$. We
identify $\Rn$ with $R^n$ by $x=\sum_ex_eX^e\leftrightarrow(x_0,\dots,x_{n-1})$
and write $C\le \Rn$ for a right $\Rn$-submodule.
\end{definition}

The choice of right ideals is forced, not conventional: see
Lemma~\ref{lem:failure}. The ring $\Bq$ is commutative, being a quotient of the
commutative ring $A_t[X]$, and there is a canonical isomorphism
\begin{equation}\label{eq:morita}
  \Rn \;\cong\; M_2(\Bq)
\end{equation}
because $X$ is central, so that a polynomial with matrix coefficients may be
read as a matrix of polynomials. Under \eqref{eq:morita}, right ideals of
$\Rn$ correspond to $\Bq$-submodules of $\Bq^2$; this is the Morita
correspondence we use repeatedly and we do not reprove it.

\begin{lemma}[basic structure]\label{lem:basic}
\begin{enumerate}[label=\textnormal{(\alph*)},leftmargin=*]
\item $A_t$ is a finite chain ring: its ideals are exactly
  $u^iA_t$, $0\le i\le t$; $|A_t|=q^t$; its maximal ideal is $uA_t$ and its
  residue field is $\F_q$.
\item $R$ is a finite Frobenius ring with $|R|=q^{4t}$,
  $\dim_{\F_q}R=4t$, $Z(R)=A_t\cdot I$, $\Jac(R)=M_2(uA_t)$, and exactly $t+1$
  two-sided ideals, namely $u^iR$.
\item $\dim_{\F_q}R^n=\Nlen=4tn$.
\end{enumerate}
\end{lemma}

\begin{proof}
(a) Every element of $A_t$ has a unique form $c_{i}u^i+\text{(higher terms)}$
with $c_i\ne0$, and then the element is $u^i$ times a unit; hence every ideal is
of the stated form and they form a chain. (b) The stated descriptions of the
centre and the Jacobson radical are immediate from the matrix-unit computation;
the two-sided ideals of $M_2(S)$ are $M_2(I)$ for ideals $I\trianglelefteq S$,
which gives $t+1$ of them. The Frobenius property is imported
\cite{Wood1999Duality}; a direct Frobenius functional is exhibited in
Lemma~\ref{lem:lambda} below. (c) is immediate from (b), since
$\dim_{\F_q}A_t=t$ and $\dim_{\F_q}R=4\dim_{\F_q}A_t$.
\end{proof}

\subsection{The Frobenius functional}\label{subsec:lambda}

Let $\lambda:A_t\to\F_q$ be the coefficient of $u^{t-1}$, and let
$\tr$ denote the matrix trace. Put
\begin{equation}\label{eq:Lambda}
  \Lambda=\lambda\circ\tr:R\longrightarrow\F_q,\qquad
  \Lambda(X)=\lambda(X_{11}+X_{22}) .
\end{equation}

\begin{lemma}\label{lem:lambda}
$\Lambda$ is symmetric ($\Lambda(XY)=\Lambda(YX)$), associative (the kernel
contains no nonzero left or right ideal), $\F_q$-linear, and non-degenerate:
the bilinear form $(X,Y)\mapsto\Lambda(XY)$ on $R$ is non-degenerate.
\end{lemma}

\begin{proof}
Symmetry: $\tr(XY)=\tr(YX)$ for matrices over any commutative ring, and $\lambda$
is $\F_q$-linear, so $\Lambda(XY)=\Lambda(YX)$.

Non-degeneracy: let $X=(X_{ij})\ne0$ and choose $(k,l)$ with $X_{kl}\ne0$;
write $X_{kl}=\sum_ac_au^a$ and let $a_0$ be the least index with $c_{a_0}\ne0$.
Take $Y$ to be the matrix whose $(l,k)$ entry is $u^{t-1-a_0}$ and whose other
entries are $0$. Then $XY$ has trace
$\sum_{i,j}X_{ij}Y_{ji}=X_{kl}u^{t-1-a_0}$, whose $u^{t-1}$-coefficient is
$c_{a_0}\ne0$. Hence $\Lambda(XY)\ne0$, and the form is non-degenerate.

Associativity (no nonzero one-sided ideal lies in $\ker\Lambda$): if
$X\ne0$, the $Y$ constructed above already gives $\Lambda(XY)\ne0$ with $Y$
arbitrary in $R$, so the right annihilator of $\Lambda$ in $R$ is zero; by
symmetry the same holds on the left.
\end{proof}

Together with the import that $R$ is Frobenius, Lemma~\ref{lem:lambda} says that
$(R,\Lambda)$ is a Frobenius algebra over $\F_q$, which is the setting we need.

\subsection{The right-module convention and the failure lemma}

\begin{lemma}[failure lemma]\label{lem:failure}
Let $S$ be a finite Frobenius ring and let
$C$ be a \textbf{left} ideal. Then the dual of $C$ with respect to the plain
product $\sum_ix_iy_i$ satisfies $|C|\cdot|C^{\perp}|=|S|^n$. For a
\textbf{right} ideal this fails in general. Concretely, for
$S=R=M_2(\F_2)$, $C=uR$ (with $t=2$, $q=2$, $n=1$, so that $u$ denotes a
nilpotent generator of $A_2$ embedded diagonally) one has
$|C|\cdot|C^{\perp}|=4\ne16=|R|$.
\end{lemma}

\begin{proof}
For a left ideal $C$, the map
$\psi_C:C\to\Hom_{\F_q}(S/C^{\perp},\F_q)$, $c\mapsto(y\mapsto\Lambda(cy))$, is
an $\F_q$-linear isomorphism by the non-degeneracy of $\Lambda$, and the
standard dimension count gives the size formula.

For right ideals the argument breaks: $\psi_C$ is no longer $\F_q$-balanced in
the way required, because the plain product is not compatible with the
right action. The stated counterexample is a direct computation: with
$R=M_2(\F_2)$ and $C=uR$ one has $|C|=4$ while the plain-product dual
$C^{\perp}=\{y:\tr(xy)=0\ \forall x\in C\}$ has size $1$, giving
$|C||C^{\perp}|=4\ne 16$.
\end{proof}

\begin{remark}\label{rem:convention}
Lemma~\ref{lem:failure} is the reason the pairing of
Definition~\ref{def:pairing} below carries the transpose: only then does the
dual of a right module come out as a right module with the correct size
(Theorem~\ref{thm:duality}). This is not a matter of notation.
\end{remark}

\subsection{Automorphisms and the anti-automorphism}

\begin{definition}\label{def:sigma}
For $s\ge0$ let $\sigma_s:\F_q\to\F_q$, $a\mapsto a^{p^s}$. Extend $\sigma_s$
to $A_t$, to $R$ and to $\Rn$ coefficient-wise, fixing $u$ and $X$. Define
\begin{equation}\label{eq:theta}
  \theta_s=\sigma_s\circ(\cdot)^{\mathsf T},
\end{equation}
the composite of the coefficient-wise $\sigma_s$ with the transpose.
\end{definition}

\begin{lemma}\label{lem:theta}
\begin{enumerate}[label=\textnormal{(\alph*)},leftmargin=*]
\item $\sigma_s\in\Aut(A_t)$ and its entrywise extension lies in $\Aut(R)$;
  $\sigma_s$ is nontrivial iff $r\nmid s$, and its order is $r/\gcd(r,s)$.
\item $\theta_s$ is an \textbf{anti}-automorphism of $R[X]$ and of $\Rn$:
  $\theta_s(ab)=\theta_s(b)\theta_s(a)$.
\item $\theta_s^2=\sigma_{2s}$; more generally $\theta_s^{2k}=\sigma_{2ks}$ and
  $\theta_s^{2k-1}=\theta_{(2k-1)s}$. In particular $\theta_s$ is an involution
  iff $r\mid 2s$.
\item $\theta_s(X^n-1)=X^n-1$, so $\theta_s$ is well defined on $\Rn$; it is
  additive, $\F_p$-linear and $\sigma_s$-semilinear over $\F_q$ (not
  $\F_q$-linear unless $r\mid s$).
\end{enumerate}
\end{lemma}

\begin{proof}
(a) $a\mapsto a^{p^s}$ is a field automorphism of $\F_q$ of order
$r/\gcd(r,s)$, and it fixes the coefficients of the relation $u^t=0$, so it
extends to $A_t$ by fixing $u$; entrywise extension preserves matrix
multiplication.

(b) For matrices over any ring, $(AB)^{\mathsf T}=B^{\mathsf T}A^{\mathsf T}$,
and $\sigma_s$ is a ring automorphism, hence
$\theta_s(ab)=\sigma_s((ab)^{\mathsf T})
=\sigma_s(b^{\mathsf T})\sigma_s(a^{\mathsf T})=\theta_s(b)\theta_s(a)$. The
same computation works in $R[X]$; note that there the index of $X$ is swapped
as well, which is why the statement is recorded separately for $R[X]$ and for
$\Rn$.

(c) $\theta_s^2(X)=\sigma_s(\sigma_s(X^{\mathsf T})^{\mathsf T})=\sigma_s^2(X)
=\sigma_{2s}(X)$; the general powers follow by induction.

(d) $\sigma_s$ fixes $X$ and the coefficients of $X^n-1$ are $0,\pm1\in\F_p$,
which $\sigma_s$ fixes. Semilinearity is immediate from $\sigma_s$ being a field
automorphism of order $>1$ in general.
\end{proof}

\begin{remark}
We use only the coefficient-Frobenius subgroup of $\Aut(A_t)$: an automorphism
of $A_t$ may move $u$, and we never use those. Everything in this paper
generalises verbatim to $\tau\in\Aut(A_t)$ with $\theta_\tau=\tau\circ(\cdot)^{\mathsf T}$
and $\theta_\tau^2=\tau^2$; we keep $\sigma_s$ for concreteness.
\end{remark}

\subsection{The CRT decomposition of \texorpdfstring{$\Bq$}{B}}%
\label{subsec:crt}

Let $X^n-1=\prod_{j}g_j^{e_j}$ be the factorisation into pairwise coprime
monic irreducibles over $\F_q$. By the Chinese remainder theorem,
\begin{equation}\label{eq:crt}
  \Bq=A_t[X]/\langle X^n-1\rangle
  \;\cong\;\prod_j \Bq_j,\qquad
  \Bq_j=A_t[X]/\langle g_j^{e_j}\rangle ,
\end{equation}
and, writing $d_j=\deg g_j$, one has $\Bq_j\cong\F_{q^{d_j}}[u,w]/\langle u^t,w^{e_j}\rangle$.
The isomorphism \eqref{eq:crt} is realised explicitly as an $\F_q$-linear ring
isomorphism $\Psi$ (with matrix and inverse), so that it composes with the Gray
map, the Morita identification and the quantum routes.

\begin{lemma}[component generation]\label{lem:components}
Let $W\le \Bq^2$. Under \eqref{eq:crt}, $W=\prod_jW_j$ with
$W_j\le\Bq_j^2$, and
\[
  \MU_{\Bq}(W)=\max_j\MU_{\Bq_j}(W_j),
\]
where $\MU_S(W)$ is the minimal number of generators of $W$ as an
$S$-module. Moreover $\Bq_j$ is a finite chain ring iff $t=1$ or $e_j=1$; in
the remaining ``wild'' case ($t\ge2$ and $e_j\ge2$) the maximal ideal of
$\Bq_j$ is $\mathfrak m_j=(u,w)$, which is not principal, and the bound
$\MU_{\Bq_j}(W_j)\le2$ fails.
\end{lemma}

\begin{proof}
Generation is componentwise because a generating set may be chosen
componentwise and the idempotents of the product allow the components to be
addressed independently; this gives the displayed formula.

For the chain-ring criterion: $\Bq_j\cong\F_{q^{d_j}}[u,w]/\langle
u^t,w^{e_j}\rangle$ is local with maximal ideal $(u,w)$ when $t\ge2$ and
$e_j\ge2$; its maximal ideal then needs two generators, so the ring is not a
chain ring. If $t=1$ the ring is $\F_{q^{d_j}}[w]/\langle w^{e_j}\rangle$, a
chain ring; if $e_j=1$ it is $\F_{q^{d_j}}[u]/\langle u^t\rangle$, a chain
ring. Finally, a finite local ring is a chain ring iff its maximal ideal is
principal, which proves the converse.

For the failure of $\MU\le2$: over the wild ring
$\F_q[u,w]/\langle u^2,w^2\rangle$ the module $\mathfrak m\oplus\mathfrak
m\subset \Bq_j^2$ requires more than two generators; explicit submodules
with $\MU=3$ are constructed in Theorem~\ref{thm:nonprincipal}.
\end{proof}

\begin{remark}[the CRT step is what opens $n=5,6$]
Earlier component routines for $\Bq$ covered only three regimes: split
semisimple, $n=p^v$ with $t=1$, and ``wild'' $n=p^v$ with $t\ge2$. The general
decomposition \eqref{eq:crt} removes that restriction, which is what makes the
families $(q,t,n)=(4,1,5)$, $(4,1,6)$ and $(8,1,3)$ reachable; they are
precisely the families from which two of the headline codes in
Section~\ref{sec:cmp} come.
\end{remark}

\subsection{The reversal operator}

\begin{definition}\label{def:rev}
Let $\rev:\Rn\to\Rn$ be the $\F_q$-linear map
$\rev\bigl(\sum_ex_eX^e\bigr)=\sum_ex_{-e}X^e$, the indices being read modulo
$n$.
\end{definition}

\begin{lemma}\label{lem:rev}
$\rev$ is a ring \textbf{automorphism} of $\Rn$ (not an anti-automorphism), and
$\rev=\id$ on $\Rn$ if and only if $n\le2$.
\end{lemma}

\begin{proof}
Multiplicativity is the computation
\[
  \rev(ab)=\sum_e\Bigl(\sum_{i+j=e}a_ib_j\Bigr)X^{-e}
  =\sum_{i,j}a_ib_jX^{-i-j}=\rev(a)\rev(b) ,
\]
which uses commutativity of $X$ with the coefficients but no commutativity of
the coefficients with each other; and $\rev$ is an automorphism because the
index map $e\mapsto -e$ is an involution. If $n\le2$ then
$-e\equiv e\pmod n$ for all $e$, so $\rev=\id$; conversely, if $n\ge3$ then
$\rev(X)=X^{-1}\ne X$.
\end{proof}

\subsection{Worked examples I: the ring and the CRT decomposition}

\begin{example}[the ring $R$ is not commutative, and $\theta_s$ is an
anti-automorphism]\label{ex:ring}
Take $q=4$, $t=2$, so that $A_2=\F_4[u]/\langle u^2\rangle$ and
$R=M_2(A_2)$. By Lemma~\ref{lem:basic}(b),
$|R|=q^{4t}=4^8=65536$ and $\dim_{\F_4}R=8$; the centre is $A_2\cdot I$ and
the Jacobson radical is $M_2(uA_2)$, so $R$ has exactly $t+1=3$ two-sided
ideals, namely $R$, $uR$ and $0$.

Let $E_{kl}$ denote the matrix units and take $s=0$, so that
$\theta_0(X)=X^{\mathsf T}$. Then
\[
  E_{12}E_{21}=E_{11},\qquad\text{so}\qquad
  \theta_0(E_{12}E_{21})=E_{11},
\]
whereas
\[
  \theta_0(E_{12})\,\theta_0(E_{21})=E_{21}E_{12}=E_{22}\ne E_{11},\qquad
  \theta_0(E_{21})\,\theta_0(E_{12})=E_{12}E_{21}=E_{11}.
\]
So $\theta_0(ab)=\theta_0(b)\theta_0(a)$ and $\theta_0$ is an
\emph{anti}-automorphism, not an automorphism. Over a commutative base ring
this distinction disappears entirely --- $\theta_s$ degenerates to $\sigma_s$
--- which is why the whole of Section~\ref{sec:cyclic} has no commutative
analogue.
\end{example}

\begin{example}[a wild CRT component]\label{ex:crt}
Take $(q,t,n)=(4,2,2)$. In characteristic $2$,
\[
  X^2-1=(X-1)^2\quad\text{over }\F_4,
\]
so the factorisation \eqref{eq:crt} has a single factor $g_1=X-1$ of degree
$d_1=1$ with multiplicity $e_1=2$. Hence
\[
  \Bq=A_2[X]/\langle (X-1)^2\rangle\;\cong\;\F_4[u,w]/\langle u^2,w^2\rangle,
\]
a local ring that is \textbf{not} a chain ring: its maximal ideal is
$\mathfrak m=(u,w)$, which needs two generators, and
\[
  \mathfrak m^2=(uw),\qquad \mathfrak m^3=0,\qquad
  \dim_{\F_4}\mathfrak m=3,\ \dim_{\F_4}\mathfrak m^2=1,\qquad
  \dim_{\F_4}\ann_{\Bq^2}(\mathfrak m)=2 .
\]
This is the ring in which $\MU_{\Bq}(W)\le2$ fails
(Lemma~\ref{lem:components}), and it is the source of all the witnesses of
Theorem~\ref{thm:nonprincipal}. The whole family has $|\Bq|=4^4=256$ and
exactly $1221$ $\Bq$-submodules of $\Bq^2$ --- the number $1221$ in the census
of Section~\ref{sec:cert} --- each corresponding to one cyclic code of length
$2$ over $M_2(A_2)$.
\end{example}
